% Generated draft: requires mathematical and source review.
\documentclass[11pt]{article}
\usepackage[margin=25mm]{geometry}
\usepackage{fontspec}
\setmainfont{Latin Modern Roman}
\usepackage{amsmath,amssymb,mathtools,unicode-math}
\setmathfont{Latin Modern Math}
\usepackage{enumitem,needspace,xurl,hyperref}
\setlength{\parindent}{0pt}
\setlength{\parskip}{5pt}
\setlength{\emergencystretch}{3em}
\newcommand{\N}{\mathbb{N}}
\newcommand{\Z}{\mathbb{Z}}
\newcommand{\Q}{\mathbb{Q}}
\newcommand{\R}{\mathbb{R}}
\newcommand{\C}{\mathbb{C}}
\newcommand{\F}{\mathbb{F}}
\newcommand{\A}{\mathbb{A}}
\newcommand{\PP}{\mathbb{P}}
\newcommand{\E}{\mathbb{E}}
\newcommand{\eps}{\varepsilon}
\newcommand{\dd}{\,\mathrm{d}}
\begin{document}
\section*{UnsolvedMath ID: 20000737}
\textbf{Source problem code:} catalogue entry
\subsection{Definitions and mathematical statement}
The original asks how $q$-analogues of the dilogarithm enter the theory, in the workshop's Nahm-sum context. It does not specify which quantum dilogarithm or what relationship is sought. The attached aid uses the compact version. For $0<q<1$ define
\[
(z;q)_\infty=\prod_{m\ge0}(1-zq^m),\qquad
\mathcal L_q(z)=\sum_{m\ge1}\frac{z^m}{m(1-q^m)}=-\log(z;q)_\infty\quad(|z|<1).
\]
Let $A\in M_r(\Q)$ be symmetric positive definite, $B\in\Q^r$, $C\in\Q$, and consider the Nahm summand
\[
w_q(\mathbf n)=\frac{q^{\frac12\mathbf n^tA\mathbf n+B^t\mathbf n+C}}{\prod_i(q)_{n_i}},\qquad \mathbf n\in\Z_{\ge0}^r.
\]
Writing $X_i=q^{n_i}$ gives the exact action identity
\[
\log w_q(\mathbf n)=\left(\tfrac12\mathbf n^tA\mathbf n+B^t\mathbf n+C\right)\log q
-\sum_i\mathcal L_q(qX_i)-r\log(q;q)_\infty.
\]
The aid states that the continuous real extension of this log-weight to $\mathbf n\in[0,\infty)^r$ is strictly concave. Its adjacent-lattice ratio is
\[
\frac{w_q(\mathbf n+e_i)}{w_q(\mathbf n)}
=\frac{q^{A_{ii}/2+B_i}\prod_jX_j^{A_{ij}}}{1-qX_i}.
\]
Thus adjacent-weight balance gives the exact finite-$q$ system
\[
1-qX_i=q^{A_{ii}/2+B_i}\prod_jX_j^{A_{ij}}.
\]
This balance equation is not asserted to be identical to the continuous derivative critical-point equation at finite $q$.

Put $q=e^{-\epsilon}$ and $u_i=-\log X_i=\epsilon n_i$. The scaled action has classical limit
\[
\Phi(u)=\sum_i\left(\frac{\pi^2}{6}-\operatorname{Li}_2(e^{-u_i})\right)-\tfrac12u^tAu,
\qquad \operatorname{Li}_2(z)=\sum_{m\ge1}\frac{z^m}{m^2}.
\]
Its unique critical point $u^*\in(0,\infty)^r$ solves
\[
1-X_i^*=\prod_j(X_j^*)^{A_{ij}},\qquad X_i^*=e^{-u_i^*},
\]
the classical Nahm balance equations used in the Bloch-group setting. The aid further states that the real-analytic local branch of the finite-$q$ balance system through $u^*$ has
\[
u'(0)=-(A+H)^{-1}(h+d),\quad
h_i=\frac{X_i^*}{1-X_i^*},\quad H=\operatorname{diag}(h_i),\quad d_i=A_{ii}/2+B_i.
\]
The derivative $u'(0)$ is with respect to $\epsilon$. The parameter $A$ controls the leading classical point, $B$ enters this first displacement, and $C$ contributes only an overall normalization of the summand. These statements give a local finite-$q$/classical dictionary, not a modularity criterion for all Nahm sums.
\subsection{Short English statement}
The compact quantum dilogarithm expresses each Nahm summand exactly. Its finite lattice balance deforms the classical dilogarithmic critical equations, and the supplied aid gives the first displacement of that balance point as $q$ tends to one.
\subsection{Sources}
\begin{itemize}
\item Local UnsolvedMath registry, permanent ID $20000737$, source code catalogue entry; canonical \texttt{aim-analytic-number-theory-notes.json}, \texttt{notes[100]}.
\item American Institute of Mathematics, ``Mock modular forms in combinatorics and arithmetic geometry,'' problem 2.6: \url{https://aimath.org/WWN/mockmodular/mockmodular.pdf}. The exact question is supplied in the registry; no original remarks or literature references are given. No PDF inspection is claimed here.
\item Attached research attempt 1: partial result, run and registry status partially solved, rigor proved, novelty confidence low, classification PARTIAL-PROGRESS. The machine-generated aid is not peer reviewed and requires independent verification before citation. No authoritative citation for the deformation lemma is supplied.
\end{itemize}
\end{document}
